\section*{第 \ref{ch:g8:contraception} 章——避孕}

\begin{solution}{exo:g8:contraception:1}
\omterm{def:g8:sexual-reproduction:gametes}{配子}产生——精液射入——\omterm{def:g8:reproductive-systems:male}{精子}游向
\omterm{def:g8:reproductive-systems:female}{输卵管}——\omterm{def:g4:animal-reproduction:sexual}{卵细胞}排出并与\omterm{def:g8:reproductive-systems:male}{精子}相遇——\omterm{def:g5:flower-to-fruit:steps}{受精}
——着床。\omterm{prop:g8:contraception:principle}{避孕}即对这一链条中的任一环节进行有意的
阻断。
\end{solution}

\begin{solution}{exo:g8:contraception:2}
\omterm{def:g8:contraception:condom}{避孕套}阻断射入环节：精液永远无法进入
伴侣体内。\omterm{def:g8:contraception:pill}{避孕药}阻断\omterm{prop:g8:reproductive-systems:cycle}{排卵}环节：不再\omterm{prop:g8:reproductive-systems:cycle}{排卵}，
也就没有可与之相遇的\omterm{def:g4:animal-reproduction:sexual}{卵细胞}。
\end{solution}

\begin{solution}{exo:g8:contraception:3}
预防经性接触传播的感染——
市面上其他\omterm{prop:g8:contraception:principle}{避孕}方法均不具此作用。
\end{solution}

\begin{solution}{exo:g8:contraception:4}
主要是推迟尚未发生的\omterm{prop:g8:reproductive-systems:cycle}{排卵}，通常在
事后数日内服用。它并非常规
\omterm{prop:g8:contraception:principle}{避孕}方法，只是最后的补救手段，不能替代
事先采用的\omterm{prop:g8:contraception:principle}{避孕}方法。
\end{solution}

\begin{solution}{exo:g8:contraception:5}
校医和家庭医生：可保密地提供解答
及开具处方，未成年人亦可。计划生育
服务中心：免费提供不带评判的咨询和\omterm{prop:g8:contraception:principle}{避孕}方法。药剂师：
紧急\omterm{def:g8:contraception:pill}{避孕药}及其服用时限。
\end{solution}

\begin{solution}{exo:g8:contraception:6}
\omterm{prop:g8:reproductive-systems:cycle}{排卵}日期无法事先得知——周期
因人而异、因月而异——而\omterm{def:g8:reproductive-systems:male}{精子}可存活
数日，等候时机。日历既无法确定\omterm{prop:g8:reproductive-systems:cycle}{排卵}时间，
也防不住等待中的\omterm{def:g8:reproductive-systems:male}{精子}。
\end{solution}

\begin{solution}{exo:g8:contraception:7}
阻断相遇：\omterm{def:g8:contraception:condom}{避孕套}（阻断射入）、宫内节育器
（令相遇和着床失败）。阻断\omterm{prop:g8:reproductive-systems:cycle}{排卵}：
\omterm{def:g8:contraception:pill}{避孕药}，以及皮下埋植剂和\omterm{prop:g8:contraception:principle}{避孕}贴片——原理相同，
只是更省心。事后补救：紧急\omterm{def:g8:contraception:pill}{避孕药}，
推迟\omterm{prop:g8:reproductive-systems:cycle}{排卵}。
\end{solution}

\begin{solution}{exo:g8:contraception:8}
\omterm{def:g8:contraception:pill}{避孕药}守住整条链条——不\omterm{prop:g8:reproductive-systems:cycle}{排卵}，因而不会\omterm{def:g5:human-reproduction:pregnancy}{怀孕}；
\omterm{def:g8:contraception:condom}{避孕套}保护伴侣——它是唯一能挡住感染的
屏障。两种不同的风险，
各需一种方法应对。
\end{solution}

\begin{solution}{exo:g8:contraception:9}
处于事件 4 所述的“暂停”状态，即
\omterm{def:g5:human-reproduction:pregnancy}{怀孕}本身的\omterm{def:g8:puberty:hormones}{激素}所造成的状态，但并未\omterm{def:g5:human-reproduction:pregnancy}{怀孕}：
只要信号持续，\omterm{def:g8:reproductive-systems:female}{卵巢}便不\omterm{prop:g8:reproductive-systems:cycle}{排卵}。
没有排出的\omterm{def:g4:animal-reproduction:sexual}{卵细胞}，就无可\omterm{def:g5:flower-to-fruit:steps}{受精}，
也无可着床。
\end{solution}

\begin{solution}{exo:g8:contraception:10}
这条链条贯穿双方身体——一方提供\omterm{def:g8:sexual-reproduction:gametes}{配子}，
另一方的周期和着床——其后果
落在双方身上。任何一方都可能阻断的环节，
双方都应共同商议：方法、记忆和执行
共同承担，而非交由一方代劳。
\end{solution}

\begin{solution}{exo:g8:contraception:11}
(a) 皮下埋植剂：环节——\omterm{prop:g8:reproductive-systems:cycle}{排卵}，通过\omterm{def:g8:puberty:hormones}{激素}持续抑制；
可靠性——名列前茅，数年内无需记忆；
感染防护——无，需与\omterm{def:g8:contraception:condom}{避孕套}配合使用；要求——
由医生放置。各项均达标，是一种正式方法。(b)
体外排精：环节——无法明确指出，因为\omterm{def:g8:reproductive-systems:male}{精子}
在最后时刻前已经游动；可靠性——正因如此而很差；
评估在问题 1 处即已终止：不是一种\omterm{prop:g8:contraception:principle}{避孕}方法。
\end{solution}

\begin{solution}{exo:g8:contraception:12}
\omterm{def:g8:contraception:pill}{避孕药}针对\cref{ch:g8:hormonal-communication}中的
系统——向\omterm{def:g8:reproductive-systems:female}{卵巢}发出\omterm{def:g8:puberty:hormones}{激素}信号；
\omterm{def:g8:contraception:condom}{避孕套}针对
\cref{ch:g8:reproductive-systems}中的示意图——在
精液射入环节设屏障；宫内节育器针对
\cref{ch:g8:from-egg-to-birth}中的第一周——相遇
与着床被阻止。每种方法都是某一章机制的
有意阻断。
\end{solution}

\begin{solution}{pb:g8:contraception:1}
\textbf{1.} 从左到右画出的链条——\omterm{def:g8:sexual-reproduction:gametes}{配子}、
射入、游动、\omterm{prop:g8:reproductive-systems:cycle}{排卵}与相遇、\omterm{def:g5:flower-to-fruit:steps}{受精}、
着床——并标出射入处的屏障、
\omterm{prop:g8:reproductive-systems:cycle}{排卵}处的\omterm{def:g8:puberty:hormones}{激素}类方法，以及相遇与着床处的
器具类方法。

\textbf{2.} 它以人体自身的\omterm{def:g8:puberty:hormones}{激素}语言，
与周期的调控链条对话：每日一剂维持
“暂停”信号，\omterm{def:g8:reproductive-systems:female}{卵巢}则听从其
所属系统的信使，不再\omterm{prop:g8:reproductive-systems:cycle}{排卵}。

\textbf{3.} 漏服若干天，信号便会中断：周期的
调控机制重新启动，\omterm{prop:g8:reproductive-systems:cycle}{排卵}可能悄然发生——
\omterm{prop:g8:reproductive-systems:cycle}{排卵}环节重新开启，且毫无预兆。

\textbf{4.} 由于各种方法都守护着生殖链条，只有
\omterm{def:g8:contraception:condom}{避孕套}能防范感染——这第二项职责
为其他方法所不具备。

\textbf{5.} \omterm{prop:g8:reproductive-systems:cycle}{排卵}日期无法预先确定——
周期长短不一，\omterm{def:g8:reproductive-systems:male}{精子}又能存活数日，
能熬过日历的预测。该应用对二者均无法测量。

\textbf{6.} 植入剂或贴片（释放保持关闭，
以月计）与宫内节育器（相遇与着床均被阻止，
以年计）——二者在设计上均无需记忆，
不必每日记得服用。

\textbf{7.} 紧急\omterm{def:g8:contraception:pill}{避孕药}；其时限——越早越好，
数日之内服用皆可；其主要机制——推迟
尚未发生的\omterm{prop:g8:reproductive-systems:cycle}{排卵}。事后则需
长期选用一种恰当的方法——紧急\omterm{def:g8:contraception:pill}{避孕药}只是刹车，并非制动系统。

\textbf{8.} “无论你用什么其他方法：只有\omterm{def:g8:contraception:condom}{避孕套}
还能同时保护你们免受感染——请把它放在
考虑之中。”

\textbf{9.} 因为能说出环节名称的使用者，
便能评估今后听到的任何方法——并看穿
每一则未提及任何环节的“传说”。中心
讲授这张图，方框才有意义。

\textbf{10.} 每日需记住：服药。植入或
放置：植入剂、节育器。仅在需要时：\omterm{def:g8:contraception:condom}{避孕套}。
看时限，去药房：紧急\omterm{def:g8:contraception:pill}{避孕药}。

\textbf{11.} 因为帮助图谱只有在安全可用时才起作用：
保密，才能让一名
十七岁的少年选择知识而非传说——
若取消保密，问题便又流落到操场上。

\textbf{12.} “中心提供的是机制图谱
及其所带来的选择——方框只是
付诸实施而已。”
\end{solution}
